\documentclass[10pt,aspectratio=169]{beamer}

\usetheme{Madrid}
\usecolortheme{default}
\setbeamertemplate{navigation symbols}{}
\setbeamerfont{frametitle}{size=\large}
\setbeamerfont{title}{size=\Large}

\usepackage{booktabs}
\usepackage{amsmath, amssymb}
\usepackage{listings}

\lstset{
  basicstyle=\ttfamily\footnotesize,
  breaklines=true,
  frame=none,
  showstringspaces=false,
}

\title{DreamLog}
\subtitle{Compression Enables Generalization in Logic Programming}
\author{Alexander Towell}
\institute{Southern Illinois University Edwardsville}
\date{April 2026}

\begin{document}

\begin{frame}
  \titlepage
\end{frame}

%%% =====================================================================
%%% Slide 1: What I want from this meeting
%%% =====================================================================
\begin{frame}{Three questions}
  \begin{enumerate}
    \setlength\itemsep{0.8em}
    \item Is this paper interesting enough for you to coauthor?
    \item Does it fit into my dissertation, or is it a side project?
    \item What would strengthening it look like before submission?
  \end{enumerate}

  \vspace{1em}
  \begin{block}{Status}
  \begin{itemize}
    \item Preprint on Zenodo (DOI \texttt{10.5281/zenodo.19490027})
    \item 9-page draft, code and experiments public on GitHub
    \item Reviewed twice (self-review via a multi-agent reviewer); open issues are known
  \end{itemize}
  \end{block}
\end{frame}

%%% =====================================================================
%%% Slide 2: The idea
%%% =====================================================================
\begin{frame}{The idea}

  \textbf{Problem.} Knowledge bases grow but do not learn. You can add
  300 facts about a family tree, and the system still cannot tell you
  whether a new person is someone's father. The facts accumulate,
  understanding does not.

  \vspace{0.8em}

  \textbf{Claim.} Compression \emph{is} learning (Solomonoff, MDL).
  The shortest KB consistent with the data generalizes best.

  \vspace{0.8em}

  \textbf{System.} DreamLog, a logic programming system with
  wake-sleep cycles inspired by DreamCoder:

  \begin{itemize}
    \item \textbf{Wake:} answer queries via SLD resolution, optionally
          call an LLM to propose rules for undefined predicates.
    \item \textbf{Sleep:} run 8 compression operations (7 symbolic,
          1 LLM-assisted), each verified against positive and
          negative queries with atomic rollback.
  \end{itemize}

  Every compression is behaviorally equivalent to the original KB on
  the verification suite. Nothing gets ``worse.''
\end{frame}

%%% =====================================================================
%%% Slide 3: Key experiment
%%% =====================================================================
\begin{frame}{The result I care most about}

  A synthetic crafting domain. 15 materials with invented names
  (\textit{lumite}, \textit{vexal}, \textit{drossite}, \ldots). The LLM
  has never seen these terms in training.

  \vspace{0.5em}
  33 checks on unseen materials / artisans / recipes, 5 runs per
  condition:

  \begin{center}
  \begin{tabular}{lccc}
    \toprule
    Condition & Recall & Precision & Rules \\
    \midrule
    No dream          & 0\%        & ---  & 0 \\
    Raw LLM prompting & 0\%        & ---  & 0 \\
    Symbolic only     & 53\%       & 59\% & 5 \\
    Full pipeline     & $64 \pm 2$\% & 64\% & 20 \\
    \bottomrule
  \end{tabular}
  \end{center}

  \vspace{0.8em}

  The raw LLM scores 0\% because the terms are invented. Symbolic
  compression alone hits 53\% without any LLM involvement, which means
  MDL-driven generalization discovers real concepts from data
  structure, not from prior knowledge. The LLM adds another 11 points
  on top.

  \vspace{0.4em}

  This is the cleanest evidence I have for the thesis.

\end{frame}

%%% =====================================================================
%%% Slide 4: Honest assessment
%%% =====================================================================
\begin{frame}{What's strong, what's missing}

  \begin{columns}[t]

  \column{0.48\textwidth}
  \textbf{Strong}
  \begin{itemize}
    \item Novel-domain test isolates compression from LLM memorization
    \item Zero false positives in either verification mode
    \item Open-world / closed-world distinction is a real
          methodological contribution
    \item Reproducible: all experiments scripted, registry tracks
          provenance
    \item Scales to 443 assertions, 78 predicates (research lab
          simulation)
    \item LLM cost under \$0.05 for all experiments combined
  \end{itemize}

  \column{0.48\textwidth}
  \textbf{Missing}
  \begin{itemize}
    \item No head-to-head ILP baseline (Popper, Metagol)
    \item No figures; a wake-sleep diagram would help
    \item Single-author empirical paper is a harder venue sell
    \item Recursive rule discovery fails (\texttt{ancestor} predicate)
    \item Scale is tens of predicates, not thousands
  \end{itemize}

  \end{columns}

\end{frame}

%%% =====================================================================
%%% Slide 5: The ask
%%% =====================================================================
\begin{frame}{The ask}

  \textbf{What coauthoring might look like.} Where I think you could
  contribute:

  \begin{itemize}
    \item Statistical rigor on the empirical claims (your ransomware
          work set a bar for how to report these)
    \item Venue strategy; a senior coauthor materially changes what's
          reachable
    \item Framing with respect to my dissertation: is this a chapter
          or a side project?
    \item Advocating for a Popper benchmark run, which is mostly
          student-time effort but needs to be scoped
  \end{itemize}

  \vspace{0.8em}

  \textbf{What I'd do.} Everything else: bug fixes, experiments,
  writing, figures, the ILP benchmark setup, revisions.

\end{frame}

%%% =====================================================================
%%% Slide 6: Open questions
%%% =====================================================================
\begin{frame}{Open questions for discussion}

  \begin{enumerate}
    \setlength\itemsep{0.8em}
    \item Dissertation fit. The listed advising scope is AI safety,
          interpretability, network analysis. DreamLog produces
          interpretable rules from unstructured KBs. Does that count?
    \item Venue. NeSy 2026 (Lisbon, September), NeurIPS workshop
          (December), or ILP 2027? Each has different bar, different
          timeline.
    \item Priority ordering. If strengthening, what goes first:
          Popper benchmark, figures, or scale?
  \end{enumerate}

\end{frame}

\end{document}
